\chapter{v1.6: Zustand, Energie und die Grenze der einfachsten Dynamik}
\label{ch:v16-state}
\section{Der gemeinsame native Operator als präziser Prüfgegenstand}
Die betrachtete einzelne Bank hat 64 fermionische und 60 bosonische Moden. In einer festen, gepinnten Vorzeichenkonvention definiert $P_A$ die durch Zeile $A$ von $W$ angegebene Zweifermionenvernichtung. Der Operator lautet
\begin{equation}
H=\Delta\sum_A b_A^\dagger b_A
+g\sum_A(b_A^\dagger P_A+P_A^\dagger b_A),\qquad \Delta>0,
\label{eq:v16-H}
\end{equation}
mit reellem $g$; eine gemeinsame Phase kann in diesem Modell absorbiert werden. Er erhält
\begin{equation}
N=N_f+2N_b.
\end{equation}
Das Modell ist ein echter Wechselwirkungsoperator und kein bloßes Wurzelbild. Seine Wahl und die physische Bedeutung seiner Moden sind dennoch Voraussetzungen. Auf dem endlichen fermionischen Fockraum sind alle $P_A$ beschränkt; die bosonischen linearen Terme haben relative Schranke null zur positiven Oszillatorenergie. Die folgende Quadratvervollständigung kann zunächst auf dem endlichen Teilchenkern und dann als geschlossene Form gelesen werden.

\section{Eine vollständige Operatoridentität statt einzelner Spektralstichproben}
Die für diese Revision unabhängig rekonstruierte Identität lautet
\begin{equation}
\boxed{\mathcal A:=\sum_A P_A^\dagger P_A
=\frac{15N_f-C_{\mathrm{Spin}(10)}-C_{SU(4)}}2.}
\label{eq:v16-casimir}
\end{equation}
Die Normierung ist $C_{\mathrm{Spin}(10)}(16)=45/4$ und $C_{SU(4)}(4)=15/4$. Der neue Prüfer baut beide Generatorfamilien aus Jordan--Wigner-Außenalgebraoperatoren auf und hebt sie auf alle $\binom{64}{2}=2016$ Zweifermionenzustände an. Für die vierfach normierten Casimirs verschwindet die vollständige Matrix
\begin{equation}
8W^\dagger W+4C_{\mathrm{Spin}(10)}+4C_{SU(4)}-120I_{2016}
\end{equation}
exakt. Alle Zwischenkomponenten sind gaußsche ganze Zahlen; die größte auftretende Magnitude ist 120 und liegt weit unter der exakten Ganzzahlschranke der benutzten Zahlendarstellung. Es wird keine kleine numerische Restnorm als Null ausgegeben.

Warum reicht das für den vollen fermionischen Fockraum? Beide Seiten sind normalgeordnete Operatoren vom Grad höchstens vier. Der Nullteilchensektor bestimmt den konstanten Anteil; der Einteilchensektor den quadratischen Anteil; der vollständige Zweiteilchensektor den quartischen Anteil. Auf null und eins verschwinden die relevanten Differenzen, auf zwei ist die gesamte Matrix null. Deshalb gilt \eqref{eq:v16-casimir} in jeder Teilchenzahl, nicht nur in numerisch ausgewählten irreduziblen Räumen.

Da die Casimirs positiv sind,
\begin{equation}
0\le\mathcal A\le\tfrac{15}{2}N_f\le480I.
\end{equation}
Die obere Norm 480 wird im vollbesetzten 64-Fermionenzustand erreicht. Daraus folgt durch Quadratvervollständigung die Stabilitätsschranke $H\ge-480g^2/\Delta$. Sie beweist Stabilität der benannten endlichen Bank, nicht eine räumliche Kontinuumstheorie.

\section{Ein exakter Zustandsvergleich auf dem gesamten Fockraum}
Erweitere \eqref{eq:v16-H} um den symmetrieverträglichen Term $\mu N$:
\begin{equation}
H_\mu=H+\mu(N_f+2N_b).
\end{equation}
Die Wahl von $\mu$ ist hier \emph{eine Gegenprobe}, kein neuer passend eingestellter TFPT-Parameter. Quadratvervollständigung liefert
\begin{align}
H_\mu={}&\Delta\sum_A\left(b_A+\frac g\Delta P_A\right)^\dagger
\left(b_A+\frac g\Delta P_A\right)
-\frac{g^2}{\Delta}\mathcal A+\mu N_f+2\mu N_b\\
\ge{}&\left(\mu-\frac{15g^2}{2\Delta}\right)N_f+2\mu N_b.
\end{align}
Bei den konkreten rationalen Werten $g=\Delta/20$, $\mu=\Delta/50$ folgt
\begin{equation}
\mu-\frac{15g^2}{2\Delta}=\frac\Delta{800},\qquad
\boxed{H_\mu\ge\frac\Delta{800}(N_f+2N_b).}
\label{eq:v16-vacuum}
\end{equation}
Der $N=0$-Raum besteht nur aus dem leeren Vakuum, dessen Energie null ist. Es ist daher der eindeutige Grundzustand, mit einer bewiesenen Anregungsuntergrenze $\Delta/800$; die wirkliche Lücke darf größer sein.

Für $\mu=0$ ist das leere Vakuum dagegen bei $g\ne0$ nicht Grundzustand. Schon der invariante Zweizustandsraum eines hellen Fermionpaares und eines Vermittlers hat Matrix
\begin{equation}
\begin{pmatrix}0&\sqrt8g\\\sqrt8g&\Delta\end{pmatrix},\qquad
E_- =\frac{\Delta-\sqrt{\Delta^2+32g^2}}2<0.
\end{equation}
Dieser Gegenvergleich benötigt weder den vollständigen Grundzustandsbeweis bei $N=64$ noch eine große Diagonalisierung. Er ist ein eigener vollständiger Satz über die fehlende Auswahl \emph{innerhalb der benannten Hamiltonfamilie}.

\begin{bild}{Was die Gegenprobe bedeutet}
Die gleichen Bauteile und die gleichen Erhaltungsgesetze erlauben eine Maschine, die im leeren Zustand ruht, und eine, deren niedrigste Energie unterhalb dieses leeren Zustands liegt. Die gemeinsame Symmetrie entscheidet zwischen ihnen nicht. Eine zusätzliche einfache Auswahlregel könnte das durchaus tun; sie muss dann ausdrücklich formuliert und geprüft werden.
\end{bild}

\section{Welche Beobachtungen die Auswahl überhaupt sehen}
Weil $[H,N]=0$, gilt für jeden neutralen Operator $O$ mit $[O,N]=0$ bei festem präpariertem Eingang
\begin{equation}
e^{itH_\mu/\hbar}Oe^{-itH_\mu/\hbar}
=e^{itH/\hbar}Oe^{-itH/\hbar}.
\end{equation}
Auch ausschließlich neutrale Mehrzeitfolgen unter demselben Eingangsvertrag unterscheiden diese Zeitentwicklungen nicht. Das behauptet \emph{nicht}, alle physikalischen Experimente seien gleich: Grundzustandspräparation, thermische Gewichte, Ladungswechsel und ein aufgeladenes Referenzsystem können unterscheiden.

Innerhalb eines bereits fest ausgewählten $N$-Sektors ist $\mu N$ nur eine additive Energiekonstante. Wer ausschließlich diesen Sektor zulässt, kann die globale Vakuumwahl damit nicht testen. Der richtige nächste Auftrag ist deshalb zweifach: zunächst den physisch erlaubten Sektor- und Präparationsvertrag angeben, dann die dafür benötigte Energiewahl aus derselben Quelle begründen. Man darf nicht nachträglich zwischen festem Sektor und globalem Grundzustand wechseln.

Die jüngste Worker-Eingabe berichtet einen eindeutigen $N=64$-Singulettgrundzustand für einen bestimmten schwachen Kopplungsbereich sowie eine vollständige endliche Vierladungsprüfung. Diese großen Resultate werden hier als Quellbefund geführt, nicht als erneut durchlaufene globale Zertifizierung. Die Schranke \eqref{eq:v16-vacuum} ist davon unabhängig frisch geprüft und widerspricht ihnen nicht: Sie betrifft $H_\mu$, nicht $H$.

\section{Warum reines Vermittlerhopping keine einzelnen Fermionen trägt}
Setzt man mehrere Banken $x$ auf einen vorgegebenen Graphen und verbindet nur ihre Vermittler durch $b_x^\dagger b_y+\mathrm{h.c.}$, so bleibt
\begin{equation}
\Pi_x=(-1)^{N_{f,x}}
\end{equation}
an jeder Bank erhalten. Der lokale Vertiz vernichtet oder erzeugt zwei Fermionen; das Hopping verändert überhaupt keine Fermionzahl. Jeder Term kommutiert deshalb mit jeder $\Pi_x$.

Ein einzelnes Fermion von Bank $x$ nach $y$ zu bewegen würde beide lokalen Paritäten ändern. Die entsprechende Übergangsamplitude ist unter dieser Dynamik exakt null. Paar- oder Vermittlertransport kann dennoch stattfinden. Ein zusätzliches $f_x^\dagger f_y$ könnte die Paritätsblockade überwinden und die globale Zahl erhalten, wäre aber ein neuer dynamischer Term. Seine Herkunft muss aus einem nativen Operationswort oder einem kontrolliert abgeleiteten effektiven Prozess folgen; ein aus ausschließlich lokal geraden Termen aufgebauter Prozess kann ihn nicht heimlich erzeugen.

Die konkrete positive Aufgabe lautet daher: Gibt es in der ursprünglichen Quelle eine räumlich verbindende Operation, die an den beiden Enden ungerade ist, oder muss die physische Anregung als anderer, zusammengesetzter Sektor identifiziert werden? Eine Antwort benötigt Matrixelemente zwischen ausdrücklich benannten Zuständen. Eine beliebige neue Kopplung wäre keine Antwort auf die Herkunftsfrage.

\section{Was die Gauß-Vereinfachung leistet und wo sie endet}
Jede Zeile von $W$ enthält acht disjunkte Fermionpaare auf 16 Moden. Der zugehörige normierte helle Zustand $\ket{\psi_A}=W^\dagger e_A/\sqrt8$ hat auf diesen Moden Einteilchenkorrelation $C=I_{16}/8$ und verschwindende anomale Zweipunktfunktion. Für ein enthaltenes Paar $(i,j)$ ist
\begin{equation}
\langle n_in_j\rangle=\tfrac18,
\qquad \langle n_in_j\rangle_{\mathrm{Wick}}=\tfrac1{64},
\qquad \Delta_{\mathrm{Wick}}=\tfrac7{64}.
\end{equation}
Die neue Rechnung prüft alle 60 Zeilen. Der Zustand ist damit in diesen physischen Fermionmoden nicht gaußsch. Eine rein quadratische Theorie mit identischen Zweipunktdaten reproduziert nicht alle nativen Antworten.

Das schließt eine einfache Beschreibung in anderen Variablen nicht aus. Die Gitter-Vertexalgebra illustriert sogar, wie freie bosonische Felder mit nichtlinearen Exponentialfeldern nichttriviale Operatorprodukte tragen. Für einen solchen Wechsel muss man jedoch den physischen Zustand, die Statistik, die Präparation und die Vierpunktfunktion mitübersetzen. Die Behauptung, eine Majorana-Kovarianz \emph{allein} löse die native Dynamik, wird durch dieses konkrete Gegenbeispiel widerlegt.
